\documentclass[11pt]{article}

\usepackage[margin=0.75in]{geometry}
\usepackage{booktabs}
\usepackage{amsmath}
\usepackage{multirow}
\usepackage{hyperref}

\title{Dual-Domain MMD Adaptation: Consolidated Final Evaluation}
\author{}
\date{\today}

\begin{document}
\maketitle

\section{System and training setup}

\subsection{Hardware}

All runs in this document were trained on a single machine: 4$\times$ NVIDIA GeForce
RTX 4090 (24\,GB each, driver 535.171.04, CUDA 12.2), an AMD EPYC 74F3 24-core
processor (96 threads, 2 sockets), and 503\,GB system RAM. Every dual-domain
adaptation and no-MMD fine-tune run used 2 of the 4 GPUs via ultralytics' native DDP
(\texttt{--device 0,1}); pretraining baselines used a single GPU per run. Software:
Ultralytics 8.4.104, PyTorch 2.5.1+cu121, Python 3.12.13.

\subsection{Common hyperparameters}

To keep the comparisons in this document meaningful, every run shares the same core
training configuration unless stated otherwise: YOLOv10n, image size 1920 (native
resolution, no upscaling), 200 max epochs with early-stopping patience 100, total
batch size 16 (8 per GPU across 2 GPUs -- three exceptions are noted per-section
below), \texttt{optimizer=auto} (resolves to SGD), \texttt{lr0=0.01},
\texttt{momentum=0.937}, \texttt{weight\_decay=0.0005}, 3 warmup epochs, and
ultralytics' standard detection augmentation defaults (mosaic, HSV jitter, flips,
etc., unchanged from stock). A few runs stopped early via the patience criterion
(within 1--11 epochs of the 200 cap) rather than diverging; this is noted where
relevant.

For the dual-domain MMD runs specifically, the following settings are shared across
all variants: RBF kernel, EMA bandwidth with momentum 0.9, features extracted after
detection layer index 10 (\texttt{mmd\_target\_layer=10}), flattened (not
global-average-pooled) before the kernel. The MMD-specific knobs that vary between
the four labeled variants used throughout this document are:
\begin{itemize}
    \item \textbf{no\_mmd}: plain fine-tune on the target domain only -- no source
    batches, no MMD term at all.
    \item \textbf{naive\_mmd}: constant \texttt{mmd\_weight = 0.8} (script default),
    no bandwidth freeze -- the un-regularized baseline.
    \item \textbf{reg\_mmd}: linearly decaying \texttt{mmd\_weight} plus a frozen
    kernel bandwidth after an early epoch, intended to prevent the bandwidth
    estimate collapsing before the target distribution has moved.
    \item \textbf{reg\_mmd\_smaller\_weight}: the same regularization scheme as
    reg\_mmd, decayed to a smaller final \texttt{mmd\_weight}.
\end{itemize}

\paragraph{A note on hyperparameter provenance.} Ultralytics' own \texttt{args.yaml}
/ checkpoint \texttt{train\_args} do not capture the MMD-specific overrides (weight,
schedule, freeze epoch) -- they are stored on a nested config object that is not part
of the serialized trainer args. The values above for \texttt{mmd\_target\_layer},
kernel, and preprocessing are certain (script defaults, never overridden across this
study). The bandwidth-freeze-epoch used per comparison group is noted inline in each
section below to the best of the project's own run-directory records; treat these as
reliable but not double-logged the way the common hyperparameters above are.

\section{Synthetic target: real\_large\_source\_syn\_small\_target}

Source = real\_large (pretrained, detached), target = syn\_small, imgsz 1920,
evaluated on the syn val set. no\_mmd/naive\_mmd/reg\_mmd ran at batch 16 on 2 GPUs;
reg\_mmd\_smaller\_weight ran at batch 12 on a single GPU (an unrelated scheduling
choice, not expected to materially affect the comparison at this batch-size range).
Bandwidth-freeze-epoch = 1 for both regularized variants.

\begin{table}[h!]
    \centering
    \setlength{\tabcolsep}{5pt}
    \begin{tabular}{lrrrr}
        \toprule
        Model & person mAP50 & person mAP50-95 & vehicle mAP50 & vehicle mAP50-95 \\
        \midrule
        no\_mmd & 0.697 & 0.361 & 0.720 & 0.503 \\
        naive\_mmd & 0.674 & 0.344 & 0.693 & 0.480 \\
        reg\_mmd & 0.671 & 0.339 & 0.690 & 0.478 \\
        reg\_mmd\_smaller\_weight & 0.677 & 0.343 & 0.693 & 0.480 \\
        \bottomrule
    \end{tabular}
    \caption{real\_large $\to$ syn\_small direction, evaluated on syn val (target domain).}
\end{table}

\noindent In this direction, no-MMD clearly beats every MMD variant on both classes --
a real gap (2--3 points mAP50, 2--2.4 points mAP50-95), not the small margins seen
elsewhere in this study. Regularization does not help here either; reg\_mmd is
marginally worse than naive\_mmd. This replicates the direction-dependent asymmetry
already noted in the project's own README results: source=real/target=synthetic
improves over baseline, while source=synthetic/target=real does not fully close the
gap -- here, MMD actively underperforms plain fine-tuning in this direction.

\section{Real target}

Source = various synthetic pretrains, target = real\_small (or real\_small\_vehicle\_only),
all evaluated on real val (target domain). Two comparisons follow.

\subsection{ClearNoon-only source vs.\ all-weather source (syn\_large)}

Both groups use bandwidth-freeze-epoch = 1 for their regularized variants (matching
hyperparameters for a fair comparison).

\begin{table}[h!]
    \centering
    \setlength{\tabcolsep}{4pt}
    \small
    \begin{tabular}{llrrrr}
        \toprule
        Source & Model & person mAP50 & person mAP50-95 & vehicle mAP50 & vehicle mAP50-95 \\
        \midrule
        \multirow{4}{*}{clearnoon} & no\_mmd & 0.634 & 0.291 & 0.873 & 0.615 \\
        & naive\_mmd & 0.623 & 0.284 & 0.866 & 0.612 \\
        & reg\_mmd & 0.630 & 0.284 & 0.870 & 0.613 \\
        & reg\_mmd\_smaller\_weight & 0.617 & 0.284 & 0.868 & 0.612 \\
        \midrule
        \multirow{4}{*}{syn\_large} & no\_mmd & 0.648 & 0.296 & 0.870 & 0.617 \\
        & naive\_mmd & 0.637 & 0.291 & 0.869 & 0.618 \\
        & reg\_mmd & 0.635 & 0.292 & 0.871 & 0.619 \\
        & reg\_mmd\_smaller\_weight & 0.639 & 0.291 & 0.871 & 0.617 \\
        \bottomrule
    \end{tabular}
    \caption{Single-weather (ClearNoon only) vs.\ all-weather (syn\_large) synthetic
    source, both sized to 6400 train images.}
\end{table}

\noindent syn\_large (all four weather conditions) beats clearnoon-only on every single
row for person, and on 3 of 4 rows for vehicle (tied/marginally behind on the 4th).
The gaps are modest (0.5--1.4 points) but consistent in direction, suggesting weather
diversity in the source domain -- not just raw image count, since both sets are the
same size -- contributes a small but real benefit to target-domain performance.

\subsection{Person diminishing under MMD, and the vehicle-only comparison}

Same source dataset (syn\_large), same regularization settings (bandwidth-freeze-epoch
= 5) for the dual-class and vehicle-only rows, allowing a direct look at what happens
to the person class as MMD strengthens, and what happens to vehicle when person is
removed from training entirely.

\begin{table}[h!]
    \centering
    \setlength{\tabcolsep}{5pt}
    \begin{tabular}{lrrrr}
        \toprule
        Model & person mAP50 & person mAP50-95 & vehicle mAP50 & vehicle mAP50-95 \\
        \midrule
        \multicolumn{5}{l}{\textit{dual-class (person + vehicle)}} \\
        no\_mmd & 0.648 & 0.296 & 0.870 & 0.617 \\
        reg\_mmd & 0.642 & 0.294 & 0.872 & 0.618 \\
        reg\_mmd\_smaller\_weight & 0.635 & 0.289 & 0.876 & 0.620 \\
        \midrule
        \multicolumn{5}{l}{\textit{vehicle-only (person removed entirely)}} \\
        no\_mmd & --- & --- & 0.885 & 0.629 \\
        reg\_mmd & --- & --- & 0.885 & 0.626 \\
        reg\_mmd\_smaller\_weight & --- & --- & 0.883 & 0.624 \\
        \bottomrule
    \end{tabular}
    \caption{Dual-class rows show person declining as MMD regularization strengthens
    (no\_mmd $\to$ reg\_mmd $\to$ reg\_mmd\_smaller\_weight), while vehicle rises over
    the same sequence. Vehicle-only rows show vehicle detection with person's
    interference removed entirely.}
\end{table}

\noindent Two things happen simultaneously as the MMD regularization gets stronger
(no\_mmd $\to$ reg\_mmd $\to$ reg\_mmd\_smaller\_weight): person mAP50 falls
monotonically (0.648 $\to$ 0.642 $\to$ 0.635) while vehicle mAP50 rises monotonically
(0.870 $\to$ 0.872 $\to$ 0.876). This is the clearest evidence in this study that MMD's
small vehicle benefit comes at person's expense, not for free. Removing person from
training entirely (vehicle-only) pushes vehicle mAP50 higher still, to 0.883--0.885
across all three variants -- confirming the person class's noisier training signal
was genuinely holding back vehicle-specific performance, consistent with the
fitness-averaging mechanism discussed earlier in this study (person's volatility
disproportionately affects which epoch gets selected as best.pt when the checkpoint
metric averages both classes equally).

\section{syn\_xlarge vs.\ syn\_large as source, to close out the study}

Same target (real\_small) and same regularization settings (mmd-weight 1.0 $\to$ 0.1,
bandwidth-freeze-epoch = 5) throughout; only the source dataset's scale changes.

\begin{table}[h!]
    \centering
    \setlength{\tabcolsep}{3pt}
    \small
    \begin{tabular}{llrrrr}
        \toprule
        Source (\# imgs) & Model & person mAP50 & person mAP50-95 & vehicle mAP50 & vehicle mAP50-95 \\
        \midrule
        \multirow{4}{*}{syn\_large (6400)} & no\_mmd & 0.648 & 0.296 & 0.870 & 0.617 \\
        & naive\_mmd & 0.637 & 0.291 & 0.869 & 0.618 \\
        & reg\_mmd & 0.642 & 0.294 & 0.872 & 0.618 \\
        & reg\_mmd\_smaller\_weight & 0.635 & 0.289 & 0.876 & 0.620 \\
        \midrule
        syn\_xlarge (10240) & \textbf{reg\_mmd} & 0.635 & 0.291 & \textbf{0.876} & \textbf{0.621} \\
        \bottomrule
    \end{tabular}
    \caption{syn\_xlarge (10240 images, all weathers) as source, same regularized MMD
    settings as syn\_large's best-performing variant, both evaluated on real val.
    syn\_xlarge's run used batch 8 on 2 GPUs (4 per GPU) rather than the batch-16
    setup used elsewhere in this document.}
\end{table}

\noindent Scaling the synthetic source from syn\_large (6400 images) to syn\_xlarge
(10240 images) gives the best vehicle result of the entire study: 0.876/0.621, matching
or edging out syn\_large's own best regularized variant while leaving person roughly
unchanged (0.635, identical to syn\_large's smaller-weight variant). The gain is modest
but consistent with the study's central premise -- more of the cheap synthetic domain's
data provides a small additional benefit on top of what the regularization already
buys, without any additional cost beyond generating more synthetic frames.

\end{document}
